%%
%% Kapitel 2 -- Grundlagen
%%
%%%%%%%%%%%%%%%%%%%%%%%%%%%%%%%%%%%%%%%%%%%%%%%%%%%%%%%%%%%%%%%%%%%%%

\chapter{Grundlagen}
\label{ch.grundlagen}

Die Kühlung der verdichteten Ladeluft erhöht deren Dichte und damit die
Ladungsmasse im Zylinder. Sie ist der konstruktiv einfachste Weg zu
höherer Leistung. Über die Ladungsmasse hinaus senkt sie das
Prozesstemperaturniveau insgesamt, was die thermische Belastung der
Bauteile verringert
\cite[S.~358]{merkerGrundlagenVerbrennungsmotorenSimulation2009}.

Wie weit die Kühlung nötig ist, hängt vom Brennverfahren ab. Beim
Dieselmotor dient sie der Ladungsmasse und der Emissionsminderung. Beim
aufgeladenen Otto- oder Gasmotor ist sie zur Vermeidung des Klopfens und
für eine günstige Schwerpunktlage der Verbrennung prinzipiell notwendig
\cite[S.~358\,f.]{merkerGrundlagenVerbrennungsmotorenSimulation2009}.
Eine niedrigere Ladelufttemperatur verschiebt die Verbrennung von der
Klopfgrenze weg. Der gewonnene Abstand ist als höheres Druckverhältnis,
höheres Verdichtungsverhältnis oder als frühere Zündung nutzbar. Am
Motor dieser Arbeit ist die Aufteilung unmittelbar sichtbar: im
Dieselbetrieb steht das Regelventil offen, geregelt wird allein im
Gasbetrieb (Abschnitt~\ref{ChATCo}).

Bild~\ref{fig.grundlagen.spirale} zeigt, wie sich eine Abweichung der
Ladelufttemperatur fortpflanzt. Eine erhöhte Temperatur nach dem Kühler
steigert die Klopfneigung, worauf die Motorsteuerung den Zündzeitpunkt
zurücknimmt. Dadurch steigt die Turbineneintrittstemperatur, und bei
Erreichen des kritischen Werts muss die Leistung zurückgenommen werden.
Zugleich sinkt mit dem späten Zündzeitpunkt das Drehmoment. Ausgeglichen
wird es über einen höheren Ladedruck, der Abgasgegendruck und
Restgasanteil anhebt, womit die Kette von vorn beginnt
\cite[S.~359\,f.]{merkerGrundlagenVerbrennungsmotorenSimulation2009}.
Für den Magerbetrieb des Großmotors entfällt die Anreicherung des
Gemischs als Schutzreaktion; an ihre Stelle tritt die Rücknahme der
Leistung. Das Werk beziffert sie: ab $55\,^\circ$C beginnt
ein Derating, das bei $70\,^\circ$C $-30\,\%$ erreicht
\cite[S.~58]{everllencese4960DFProject2026}.

\begin{figure}[htb]
  \centering
  \begin{tikzpicture}[
      knoten/.style={
        draw=abbGrau, line width=0.5pt, rounded corners=2pt,
        fill=white, align=center, inner sep=4pt,
        text width=26mm, minimum height=11mm,
        font=\sffamily\footnotesize, text=abbGrau},
      zustand/.style={
        knoten, draw=abbBlau, line width=1.1pt, text=abbBlau,
        fill=abbBlau!6},
      folge/.style={
        knoten, draw=abbHell, text=abbHell},
      pfeil/.style={
        -{Stealth[length=2.2mm,width=1.6mm]}, draw=abbGrau, line width=0.7pt},
      pfeilhell/.style={
        -{Stealth[length=2.2mm,width=1.6mm]}, draw=abbHell, line width=0.6pt}
    ]
    \node[zustand] (temp)   at (  90:38mm) {Ladelufttemperatur\\steigt};
    \node[knoten]  (klopf)  at (  18:38mm) {Klopfneigung\\steigt};
    \node[knoten]  (zuend)  at ( -54:38mm) {Zündung wird nach\\spät verstellt};
    \node[knoten]  (moment) at (-126:38mm) {Drehmoment\\sinkt};
    \node[knoten]  (druck)  at ( 162:38mm) {Ladedruck und\\Restgasanteil steigen};
    \node[folge]   (turbo)  at (64mm,-31mm) {Turbineneintritts-\\temperatur steigt};
    \node[folge]   (derate) at (64mm,-53mm) {Leistung wird\\zurückgenommen};
    \draw[-{Stealth[length=2.0mm,width=1.4mm]}, draw=abbHell!55,
          line width=0.6pt]
      (135:17mm) arc[start angle=135, end angle=-155, radius=17mm];
    \draw[pfeil] (temp)   to[bend right=20] (klopf);
    \draw[pfeil] (klopf)  to[bend right=20] (zuend);
    \draw[pfeil] (zuend)  to[bend right=20] (moment);
    \draw[pfeil] (moment) to[bend right=20] (druck);
    \draw[pfeil] (druck)  to[bend right=20] (temp);
    \draw[pfeilhell] (zuend) -- (turbo);
    \draw[pfeilhell] (turbo) -- (derate);
  \end{tikzpicture}
  \caption{Wirkungskette einer erhöhten Ladelufttemperatur im
    Gasbetrieb. Nach \cite[Bild~8-39,
    S.~359]{merkerGrundlagenVerbrennungsmotorenSimulation2009},
    angepasst für den Magerbetrieb.}
  \label{fig.grundlagen.spirale}
\end{figure}

Aus dieser Kette allein folgte, dass eine möglichst kalte Ladeluft
günstig ist. Nach unten begrenzen jedoch zwei Vorgänge, die mit dem
Klopfen nichts zu tun haben. Unterhalb des Taupunkts fällt im Kühler
Wasser aus und gelangt in den Zylinder; der Projektguide rechnet für
$9\,000\,$kW bei $35\,^\circ$C und $80\,\%$ relativer Feuchte mit
$248\,$kg/h Kondensat
\cite[S.~333--334]{everllencese4960DFProject2026}. Unter
Tropenbedingungen und $3\,$bar Ladedruck liegt der Taupunkt bei rund
$51{,}5\,^\circ$C
\cite[S.~266\,f.]{pucherLadeluftkuehlungUndLadeluftkuehler2012}. Bei
niedriger Last kommt die Verbrennungsstabilität hinzu. Eine kalte Ladung
verlängert den Zündverzug, und im mageren Gasbetrieb steigt der Anteil
unverbrannten Methans im Abgas
\cite[S.~811--850]{almbauerThermomanagement2019}. Das Werk zieht daraus
eine Freigabegrenze: unter $10\,\%$ Last ist Gasbetrieb nur zulässig,
wenn die Ladelufttemperatur mindestens $50\,^\circ$C erreicht
\cite[S.~79]{everllencese4960DFProject2026}. In der Gegenrichtung wächst
mit der Ladelufttemperatur die Stickoxidbildung
\cite[S.~277]{merkerGrundlagenVerbrennungsmotorenSimulation2009}. Der
Sollwert liegt deshalb im Inneren eines Bandes und nicht an dessen Rand.

Die Werksdokumentation fasst beide Richtungen in Grenzwerte.
Für den 49/60DF ist am Zylindereintritt ein Band von $50$ bis
$60\,^\circ$C zulässig \cite[S.~162]{everllencese4960DFProject2026}.
Es beschreibt das Regelziel im Gasbetrieb; im Dieselbetrieb arbeitet
keine Sollwertregelung, und die Temperatur läuft mit der Last
(Abschnitt~\ref{ChATCo}). Die Bandkante $50\,^\circ$C ist die absolute
Betriebsgrenze, der Taupunkt des Tropen-Auslegungsfalls liegt darüber.
Das erklärt, warum der Sollwert im oberen Teil des Bandes geführt wird
(Abschnitt~\ref{ChATCo}). Die genannten Zahlen gelten für den 49/60DF;
Aufbau und Wirkmechanismen übertragen sich auf die Feldbaureihen, die
Bänder nicht (Kapitel~\ref{ch.datenbasis}). Die Untergrenze
kehrt in Kapitel~\ref{ch.ergebnisse} als weiche Ausgangsbeschränkung des
modellprädiktiven Reglers wieder, die Obergrenze begründet das enge
Sollwertband.
% PRÜFEN (Lektor C5): ch4 zitiert das Band mit "PG 49/60 Marine S. 133
% Tab. 73" — das ist die Diesel-Ausgabe (everllencese4960Project2025);
% hier die DF-Marine-Ausgabe S. 162 Tab. 90. Beide können stimmen, ch4
% sollte denselben Key nutzen.
% PRÜFEN: 51/60DF-Feldflotte hat ein anderes Band (PG 51/60DF-M S. 117
% Tab. 74: 40-70 °C), gehört nach Kap. 4/6.4; Key fehlt noch in Masterarbeit.bib.


Wie tief gekühlt werden kann, begrenzt das Kühlmedium. Die Ladeluft
erreicht bestenfalls dessen Eintrittstemperatur; darüber hinaus
bestimmen Übertragungsfläche und Kühlmittelmassenstrom das Ergebnis
\cite[S.~360]{merkerGrundlagenVerbrennungsmotorenSimulation2009}. Für
Schiffs- und Stationäranlagen steht mit dem Seewasser ein kaltes Medium
in großer Menge bereit, während Fahrzeuge auf Umgebungsluft oder
Motorkühlwasser angewiesen sind. Daraus folgt der wassergekühlte Aufbau
des hier betrachteten Ladeluftkühlers. Die Differenz zwischen
Verdichteraustritt und Kühlwassereintritt bleibt für die Modellbildung
maßgeblich; Abschnitt~\ref{ssec.grundlagen.luftpfad.kennzahlen} führt
sie als Kühlerspanne ein.

Aufladung und Ladeluftkühlung wirken auf dieselbe Größe, die Dichte der
Zylinderladung. Der Kühler gehört deshalb zum Aufladesystem, dessen
Aufbau der folgende Abschnitt beschreibt.


\section{Ladeluftsystem}

Bild~\ref{fig.grundlagen.ladeluftsystem} zeigt den Aufbau des Ladeluftsystems eines zweistufig aufgeladenen Motors.

\begin{figure}[htb]
  \centering
  \includegraphics[width=0.85\textwidth]{ladeluftsystem_zweistufig_51_60df}
  \caption{Ladeluftsystem bei zweistufiger Aufladung mit LP- und HP-Verdichter, rote Systemgrenze kennzeichnet einstufige Aufladung}
  \label{fig.grundlagen.ladeluftsystem}
\end{figure}

Beim 49/60DF wird die Ansaugluft zunächst in der Niederdruckstufe (LP)
verdichtet und im LP-Ladeluftkühler gekühlt. Es wird nicht auf die Endtemperatur gekühlt, sondern nur ausreichend, um Kondensat zu vermeiden und so den folgenden HP-Verdichter vor Beschädigung durch Tröpfchen zu schützen. In der anschließenden Hochdruckstufe (HP) wird
die Ladeluft ein zweites Mal verdichtet und im HP-Kühler auf die Endtemperatur
gebracht.

Die angesaugte Luft wird im Verdichter des Turboladers auf Ladedruck gebracht,
um die Luftdichte und damit die Sauerstoffmasse im Zylinder zu erhöhen. Aufgrund der
näherungsweise isentropen Verdichtung steigt die Ladelufttemperatur dabei
deutlich an und muss im Ladeluftkühler wieder abgesenkt werden. Das Abgas aus den
Zylindern treibt die Turbine im Abgasstrang an. Die genaue Turbinenleistung wird über die verstellbare Turbinenfläche eingestellt. Turbine und Verdichter sitzen auf einer gemeinsamen Welle; über sie wird der Verdichter angetrieben.

Der Katalysator zur selektiven katalytischen Reduktion (SCR) wirkt, anders
als eine Abgasrückführung, nicht auf die Ladelufttemperatur. Er und die im
Betrieb überwiegend inaktiven Bypässe bleiben im Weiteren unberücksichtigt.

Die zweistufige Aufladung lässt sich durch die rote Systemgrenze auf das System einer einstufigen Aufladung reduzieren. Auch aktuell wird die eigentliche Ladelufttemperatur nur über die HP-Stufe eingestellt.

\subsection{Kühlkreise}

\label{ssec.grundlagen.luftpfad.cabv}

Der Ladeluftkühler ist zweistufig aufgebaut. Die Ladeluft durchströmt zuerst das
Bündel des Hochtemperatur-Kühlkreises (HT, Kühlmedium HTCW) und anschließend das
Bündel des Niedertemperatur-Kühlkreises (LT, Kühlmedium LTCW), das die
Austrittstemperatur bestimmt. Bild~\ref{fig.grundlagen.kuelkreise} zeigt beide
Kreise mit ihren Verbrauchern.  

\emph{Stufe} bezeichnet im Folgenden stets die Aufladestufe (LP/HP),
\emph{Bündel} die HT- beziehungsweise LT-Seite innerhalb eines Kühlers.
% (Copilot 01.09.: Konventionssatz von Kommentar auf Fließtext gehoben,
% die Sachprüfung vom 27.08. hatte genau diese Konvention bestätigt.)

\begin{figure}[htb]
  \centering
  \includegraphics[width=0.85\textwidth]{kuelkreise}
  \caption{HT- und LT-Kühlkreise}
  \label{fig.grundlagen.kuelkreise}
\end{figure}

% ===== NEUFASSUNG Copilot 01.09.2026 (Fragmente Werksangabe/Sollwert
% ===== integriert; Altfassung entfernt 03.09., s. ch2.tex.bak_pre_4d0309) =====
Der HT-Kreis hält primär Zylinderköpfe und Laufbuchsen auf
Betriebstemperatur. Seine Vorlauftemperatur wird über ein Dreiwegeventil
eingeregelt, das den Durchfluss durch den HT/LT-Kühler einstellt;
dieser ist mit dem kälteren LT-Kreis verbunden und führt die Wärme
dorthin ab. Die Pumpe des HT-Kreises wird vom Motor angetrieben, der
Volumenstrom ist damit an die Drehzahl gekoppelt. Nach der Pumpe
durchströmt das HTCW das HT-Bündel der Ladeluftkühlung und anschließend
die Zylinderkühlung. Für den Start nach Stillstand liegt ein Vorwärmer
im Kreis. Wegen des Vorwärmers, der großen Wassermasse und der Lage des
HT-Bündels vor dem Motor als größtem Wärmeerzeuger schwankt die
Temperatur dieses Kreises im Betrieb wenig; wenn, dann kurzfristig als
Folge einer Lasterhöhung \cite{everllencese4960DFProject2025}.

Die Vorlauftemperatur des LT-Kreises wird ebenfalls über ein Ventil
geregelt, das den Durchfluss durch einen Kühler einstellt. Dieser
Wärmeübertrager wird von einer externen Kühlwasserquelle durchflossen,
in marinen Anwendungen vom Seewasser. Die niedrigste erreichbare
Kühlwassertemperatur hängt damit von den Umgebungsbedingungen ab
(Werksangabe $32\,^\circ$C, Betriebsband $25$ bis $38\,^\circ$C)
\cite{everllencese4960DFProject2025}. Das Temperaturniveau folgt der
Umgebung; Last und Ventilhub verschieben es leicht
(Abschnitt~\ref{sec.erg.ol}). Nach der Kühlung durchfließt das LTCW das
LT-Bündel der Ladeluftkühlung und anschließend weitere Kühlstellen,
darunter den HT/LT-Kühler, den Schmiermittelkühler und den Düsenkühler.



\subsection{Ladelufttemperaturregelung}
\label{ChATCo}
% ===== NEUFASSUNG Copilot 01.09.2026 (Altfassungen entfernt 03.09.,
% ===== s. ch2.tex.bak_pre_4d0309) =====
Die Ladelufttemperatur wird in der Serie durch die
Charge-Air-Temperature-Control (ChATCo) geregelt. Ihr Sollwert ist
keine feste Zahl: Die Werksdokumentation führt ihn abhängig vom
Betriebspunkt, von der Kraftstoffart, von Umgebungs- und
Medientemperaturen und vom Optimierungsziel
\cite[S.~51]{everllencese4960DFProject2026}; unterhalb einer
Ansauglufttemperatur-Schwelle bleibt er konstant, darüber wird er
mitgeführt \cite[S.~376]{everllencese4960DFProject2026}. Für die
Regleruntersuchungen dieser Arbeit sind $52\,^\circ$C festgelegt
(Kapitel~\ref{ch.datenbasis}); der Wert liegt innerhalb des zu
Kapitelbeginn genannten Bandes. Bild~\ref{fig.grundlagen.chatco} zeigt
die Struktur der Serienregelung.

\begin{figure}[htb]
  \centering
  \includegraphics[width=0.85\textwidth]{ChAtco}
  \caption{Aktuell verbaute Ladelufttemperaturregelung}
  \label{fig.grundlagen.chatco}
\end{figure}

Der Wasserdurchsatz durch das LT-Bündel des HP-Kühlers wird durch ein
elektrisch angesteuertes Dreiwegeventil eingestellt
\cite{everllenceseTemperaturregelventilMitAnbau2025}.
% ===== ERGÄNZT Copilot 01.09.2026 (Ventil-/Reglerstruktur, Fassung n. Lektor) =====
Die Anlage besitzt zwei baugleiche Mischventile mit linearer Kennlinie,
eines je Kühlstufe (nach dem LP- und nach dem HP-Kühler), jeweils mit
Positions-Sollwert und Positionsrückmeldung
\cite[S.~274--275]{everllencese4960DFProject2026}; maßgeblich für die
Temperatur vor Zylinder ist das HP-seitige Ventil, dessen Stellung im
Folgenden mit $\valveCA$ bezeichnet wird. Je Ventil arbeitet
ein Regler mit PI-Charakteristik und Zusatzeingängen, darunter die
Umgebungslufttemperatur (Messstelle 1TE6000)
\cite[S.~263]{everllencese4960DFProject2026}; der
Temperatursensor sitzt stromab des Ventils
\cite[S.~237]{everllencese4960DFProject2025}. Die Rohrstrecke zwischen
Mischventil und Kühlereintritt gilt in der Werksauslegung als
totzeitkritisch \cite[S.~381]{everllencese4960DFProject2026}; die
endliche Stellgeschwindigkeit des Antriebs wirkt auf die Regelung als
Ratenbegrenzung, eine Totzeit entsteht am Antrieb nur beim
Richtungswechsel (Kapitel~\ref{ch.datenbasis}). Aktiv angesteuert
wird das Ventil nur im Gasbetrieb; dort wird die Ladelufttemperatur auf
den Sollwert eingeregelt. Der Regler besteht aus einer adaptiven
Vorsteuerung und einem PI-Anteil. Die Vorsteuerung basiert auf einem
Kennfeld über der relativen Motorleistung und der Differenz zwischen
Ladeluft-Sollwert und LTCW-Temperatur; bei stationärer Last wird das
Kennfeld so nachgeführt, dass der PI-Anteil gegen null geht. Der Regler
arbeitet im Takt von $500\,$ms, der Stellbereich ist auf $0$ bis
$100\,\%$ begrenzt, und ein Anti-Windup begrenzt den Integrator an den
Stellgrenzen. Zwischen Reglerausgang und Ventilposition liegt eine
hinterlegte Linearisierungskurve, die die nichtlineare Kennlinie des
Stellkanals ausgleichen soll; die Ventilkennlinie selbst ist je Pfad
nahezu linear, die S-Form der Streckenverstärkung entsteht in Hydraulik
und Kühler (gemessen: Abschnitt~\ref{ssec.erg.ol.kennlinie}).
% Quelle der Kennlinienaussage: Clorius G3FM je Pfad "almost linear",
% RPC-Positionierer linear; die S-Form sitzt in Hydraulik und Kühlerphysik.
% Hardware-Details und die Stellraten gehören nach Kap. 4, s.
% notizen\CH2_VORSCHLAEGE_OFFEN_0309.md.


In der Steuerung sind mehrere Betriebsarten hinterlegt. Im
Dieselbetrieb steht das Ventil $100\,\%$ offen, es wird maximal
gekühlt. HT/LT-Switchover-Preheating (der LT-Kühler wird vom LT- an
den HT-Kreis gehängt) und LT-Preheating (Ventil zu) sind weitere
Betriebsarten mit eigener Dynamik. Sie werden im Weiteren nicht
betrachtet, müssen aber bekannt sein, damit Datenabschnitte
verschiedener Betriebsarten nicht vermischt werden.

Da diese Regelung bei allen Motoren aktiv ist, liegt in allen historischen Messdaten bei der Identifikation ein Closed-Loop-Identifiaktion Problem vor. Dies wird in \ref{closed_loop_identifkation} erläutert.


%%%%%%%%%%%%%%%%%%%%%%%%%%%%-------------------------------------------
% Wärmeübertragung im Ladeluftkühler
%%%%%%%%%%%%%%%%%%%%%%%%%%%-------------------------------------------


\section{Wärmeübertragung im Ladeluftkühler}
\label{ssec.grundlagen.luftpfad.waermeuebertragung}
\label{ssec.grundlagen.luftpfad.kennzahlen}
\label{ssec.grundlagen.terme.literatur}

% COPILOT-STICHPUNKTE BEGIN Herleitungsweg-20260922
% Copilot: auf Thomas' Wunsch vom 22.09.2026 auf die einfache Grundherleitung verdichtet.
% Gepruefter Modellanschluss, keine neue Modellwahl:
% 05_Matlab_TB20/21_Modellstudien/02_Koordination_Modellaufbau/
% MODELLAUSWAHL_20260922.md, Z. 10--18, 31--36 und 58--64:
% vorgeschlagene Schreibrollen EB_E4_STANDARD/S1G_ORIGINAL und E6r/UE2.
% EB21: Erklaerbare_Bibliothek_S1G/eb21_features.m, Z. 23--35:
% dH=T_col_in-T_HTCW; d3=(T_HTCW-T_LTCW)+cFix*dH;
% Tstat=T_col_in-PH*dH-P3*d3; gamma=1-exp(-rho*r*m).
% EB_E4: rate_mode R1, also r=1. cFix ist nicht exakt 1-PH.
% Die exakte Serienidentitaet unten motiviert diese Form, beweist keine
% physikalische Identifikation der voneinander unabhaengigen Fitkoeffizienten.
\par\smallskip
\noindent\textbf{Copilot: Einfache Grundherleitung für die spätere Modellbildung}
\par\smallskip
\begin{enumerate}
  \item \textbf{Wärmeabgabe der Luft.}
    Der Luftmassenstrom $\dot m_L$ gibt die pro Zeit durchströmende
    Luftmasse an. Die spezifische Wärmekapazität $c_{p,L}$ beschreibt,
    wie stark sich die Enthalpie je Luftmasseneinheit mit der Temperatur
    ändert. Bei konstantem $c_{p,L}$ ergibt sich im stationären Zustand
    \begin{equation}
      \label{eq.grundlagen.basic.luftleistung}
      \dot Q_L=\dot m_L c_{p,L}(\Tin-\Tout).
    \end{equation}
    $\dot Q_L$: von der Luft abgegebene Wärmeleistung in $\mathrm{W}$.
    Bei gleicher Abkühlung $\Tin-\Tout$ muss für einen größeren
    Luftmassenstrom mehr Wärme abgeführt werden. Bei gleicher
    Wärmeleistung fällt die Abkühlung entsprechend kleiner aus.
    Voraussetzungen: gleicher Luftmassenstrom am Ein- und Austritt,
    kein Phasenwechsel, vernachlässigbare Änderungen der kinetischen
    und potenziellen Energie. Grundlage: Enthalpieänderung und erster
    Hauptsatz nach \cite[S.~821, Gl.~31.5--31.6]{almbauerThermomanagement2019}.

  \item \textbf{Wärmeübertragung zum Kühlwasser.}
    Im Kühlfall fließt Wärme von der Luft zur Wand, durch die Wand und weiter
    zum Wasser. Zusammengefasst gilt
    \begin{equation}
      \label{eq.grundlagen.basic.waermedurchgang}
      \dot Q_{L\to W}=kA\,\Delta T_{L-W,\mathrm{mittel}}.
    \end{equation}
    $kA$: thermischer Leitwert in $\mathrm{W/K}$;
    $\Delta T_{L-W,\mathrm{mittel}}$: zur räumlichen Beschreibung
    passende mittlere Temperaturdifferenz zwischen Luft und Wasser.
    Sie ist von der luftseitigen Abkühlung $\Tin-\Tout$ zu unterscheiden.
    Die hintereinandergeschalteten Wärmeübergänge und die Wandleitung
    bestimmen den Leitwert
    \cite[S.~256, Gl.~12.3--12.7]{pucherLadeluftkuehlungUndLadeluftkuehler2012}.
    Luft- und Wasserdurchsatz beeinflussen die Wärmeübergänge und damit
    $kA$ \cite[S.~838--839, Gl.~31.59--31.62]{almbauerThermomanagement2019}.
    Die Ventilstellung wirkt über den Wasserdurchsatz. Ihre Wirkung
    auf die Kühlung muss deshalb nicht proportional zum Stellwert sein.

  \item \textbf{Temperaturdifferenzen der HT- und LT-Stufe.}
    Für die vereinfachte Reihenschaltung zunächst die Abkühlung in der
    HT-Stufe beschreiben, anschließend die weitere Abkühlung in der
    LT-Stufe:
    \begin{equation}
      \label{eq.grundlagen.basic.stufen}
      \begin{aligned}
        T_{\mathrm{zw}}&=\Tin-\varepsilon_{\mathrm{HT}}
                                   (\Tin-\THT),\\
        \Tstat&=T_{\mathrm{zw}}-\varepsilon_{\mathrm{LT}}
                                   (T_{\mathrm{zw}}-\TLT).
      \end{aligned}
    \end{equation}
    $T_{\mathrm{zw}}$: berechnete, ungemessene Lufttemperatur zwischen
    den Stufen. $\Tstat$: berechnete stationäre Motoreintrittstemperatur
    bei festgehaltenen Betriebsgrößen. Die dimensionslosen
    Temperaturänderungsgrade $\varepsilon_{\mathrm{HT}}$ und
    $\varepsilon_{\mathrm{LT}}$ beschreiben den jeweils abgebauten
    Anteil der Temperaturdifferenz zur Wassereintrittstemperatur
    \cite[S.~257, Gl.~12.8]{pucherLadeluftkuehlungUndLadeluftkuehler2012}.
    Ihre Abhängigkeit von Durchsatz und Ventilstellung folgt aus dem
    Wärmedurchgang; die konkrete Funktion bleibt zu bestimmen.
    Bei wärmerem Kühlwasser kann eine Stufe die Luft auch erwärmen;
    die entsprechende Luft-Wasser-Temperaturdifferenz wird dann negativ.
    Durch Einsetzen der ersten Stufe ergibt sich für die zweite
    \begin{equation}
      \label{eq.grundlagen.basic.ltspanne}
      T_{\mathrm{zw}}-\TLT
       =(\THT-\TLT)+(1-\varepsilon_{\mathrm{HT}})(\Tin-\THT).
    \end{equation}
    Eigene Algebra: Die LT-Stufe sieht die nach der HT-Stufe verbleibende
    Temperaturdifferenz. Damit sind genau die Kombinationen
    $\Tin-\THT$ und $\THT-\TLT$ begründet, auf die sich die spätere
    Modellierung beziehen kann. Die angenommene Stufenfolge und
    Durchströmung sind Teil der Systemvereinfachung.

  \item \textbf{Produkte aus Betriebspunktgrößen und Temperaturdifferenzen.}
    Gl.~(\ref{eq.grundlagen.basic.luftleistung}) enthält unmittelbar
    ein solches Produkt. In den Stufengleichungen steht stattdessen
    ein durchsatzabhängiger Temperaturänderungsgrad vor einer
    Temperaturdifferenz. Für seine datengetriebene Näherung kann
    beispielsweise gelten
    \begin{equation}
      \label{eq.grundlagen.basic.produkte}
      \begin{aligned}
        \varepsilon(\mrel,\valveCA)\,\Delta T
        &\approx
        \bigl[a_0+a_1\mrel+a_2g(\valveCA)\bigr]\Delta T\\
        &=a_0\Delta T+a_1\mrel\,\Delta T
                         +a_2g(\valveCA)\,\Delta T.
      \end{aligned}
    \end{equation}
    $\mrel$: dimensionsloser Luftmassenstromproxy aus den verfügbaren
    Messgrößen; $g(\valveCA)$: dimensionslose Ventilfunktion;
    $\Delta T$: jeweils die in Gl.~(\ref{eq.grundlagen.basic.stufen})
    verwendete Temperaturdifferenz. Die Koeffizienten $a_i$ werden
    aus Daten bestimmt. Wechselwirkungen können ergänzend durch
    $\mrel\,g(\valveCA)\,\Delta T$ beschrieben werden.
    Diese Zerlegung ist eine angebotene Modellnäherung. Die Physik
    begründet die Produktstruktur, aber weder bestimmte Koeffizienten
    noch eine zwingend lineare Massenstromabhängigkeit.
    Konstante Faktoren lassen sich bei geeigneter Modellform in die
    Koeffizienten aufnehmen. Die Wahl und Berechnung des Proxys wird
    unten begründet; sie liefert keine kalibrierte Massenstrommessung.

  \item \textbf{Verzögerte Temperaturantwort.}
    Wand und Fluide speichern thermische Energie. Solange sich diese
    Energie ändert, sind zugeführte und abgeführte Wärmeleistung
    verschieden. Eine zusammengefasste Speicherkapazität motiviert
    eine verzögerte Temperaturantwort
    \cite[S.~830--831, Gl.~31.36--31.38]{almbauerThermomanagement2019},
    \cite[S.~230, Gl.~4.6.14--4.6.16]{isermannEngineModelingControl2014}.
    Eine einfache eigene Modellannahme lautet
    \begin{equation}
      \label{eq.grundlagen.basic.dynamik}
      \frac{\mathrm d\Tout}{\mathrm dt}
       =\lambda(\mrel,\valveCA)(\Tstat-\Tout),
       \qquad \lambda>0.
    \end{equation}
    Die Temperatur nähert sich dem momentanen stationären Zielwert;
    $\lambda$ bestimmt die Geschwindigkeit, bei festgehaltenen
    Eingängen ist $\tau=1/\lambda$ die Zeitkonstante.
    Ein größerer Luftmassenstrom verändert den Enthalpietransport
    und kann deshalb auch diese Rate beeinflussen.
    $\lambda=c_\lambda\mrel$ ist eine mögliche einfache Näherung.
    Ein fehlender Ventilterm in der Rate ist eine Modellvereinfachung;
    die Physik schließt seinen Einfluss nicht aus.

  \item \textbf{Anschluss an die Modellierung.}
    Die Herleitung liefert als gemeinsame Bausteine
    Temperaturdifferenzen, betriebspunktabhängige Faktoren, deren
    Produkte sowie einen thermischen Speicher.
    Kapitel~\ref{ch.methodik} kann für die Kandidatenbibliothek auf
    Gl.~(\ref{eq.grundlagen.basic.produkte}), für den stationären
    Zielwert auf Gl.~(\ref{eq.grundlagen.basic.stufen}) und für die
    Dynamik auf Gl.~(\ref{eq.grundlagen.basic.dynamik}) verweisen.
    \sindyc und \edmdc können diese Bausteine unterschiedlich
    verwenden; eine lokale Linearisierung begründet eine lineare
    Vergleichsstruktur. Die nach Datenanpassung entstehenden
    Gewichte sind nicht automatisch identifizierte physikalische
    Stufenwirkungsgrade. Auswahl und Prüfung der konkreten Modelle
    gehören nach Kapitel~\ref{ch.methodik} und \ref{ch.ergebnisse}.
\end{enumerate}
% COPILOT-STICHPUNKTE END Herleitungsweg-20260922

% COPILOT-STICHPUNKTE BEGIN ch2-08
\par\smallskip
\noindent\textbf{Copilot: Stichpunkte zur Ausformulierung (16.09.2026)}
\par\smallskip

% ---- Absatz 1: Einordnung ----
\begin{itemize}
  \item Herleitung ist Vorwissen, kein gerechnetes Prozessmodell; sie legt fest,
    welche Termfamilien in Kapitel~\ref{ch.methodik} zur Wahl stehen. Drei Ebenen
    trennen: Literaturgleichung, eigene Reduktion, datengetriebene Auswahl
    (Kapitel~\ref{ch.methodik} und \ref{ch.ergebnisse}). Ein physikalisch
    plausibler Kandidat ist weder nachgewiesener Mechanismus noch automatisch
    identifizierbar.
  \item Rahmen: Mittelwertmodellierung des Luftpfads
    \cite{guzzellaIntroductionModelingControl2010,
    isermannEngineModelingControl2014}, auf Schiffsmotoren übertragen von
    \cite{theotokatosCycleMeanValue2010} und
    \cite{baldiDevelopmentCombinedMean2015}.
  \item Systemgrenze: HP-Kühler mit HT- und LT-Bündel
    (Abschnitt~\ref{ssec.grundlagen.luftpfad.cabv}). $\Tin$ wird gemessen; er
    steigt mit Ladedruck und Last \cite[S.~769--780]{baarAufladeverfahren2019},
    die Verdichterdynamik ist im Zeitbereich von $\Tout$ rasch ausgeregelt
    \cite{theotokatosCycleMeanValue2010}. Deshalb keine Verdichter- oder
    Turbinensimulation; ob $\Tin$ Eingang oder eigener Zustand ist, entscheidet
    Kapitel~\ref{ch.methodik} am Transientenfehler.
\end{itemize}

% ---- Absatz 2: stationaere Bilanz ----
% COPILOT-STICHPUNKTE BEGIN Systembilanz-20260922
\begin{itemize}
  \item \textbf{Copilot: Systemgrenze der folgenden Bilanz.}
    Für beide HT-/LT-Stufen zusammen gilt stationär ohne Außenverlust
    $\dot Q_L=\dot Q_{\mathrm{HT}}+\dot Q_{\mathrm{LT}}$.
    Die Gleichsetzung in Gl.~(\ref{eq.energiebilanz}) mit nur dem
    LT-Wasserstrom ist eine einstufige Ersatzbeschreibung.
    Das Weglassen des unabhängigen HT-Beitrags ist eine zusätzliche
    Modellvereinfachung; Grundlage sind die getrennten Stufenbilanzen.
\end{itemize}
% COPILOT-STICHPUNKTE END Systembilanz-20260922
\begin{itemize}
  \item Zuerst die übertragene Wärmeleistung, Gl.~(\ref{eq.energiebilanz}):
    Luftseite, Wärmedurchgang und Wasserseite sind gleich. Annahmen: stationär,
    keine Kondensation, adiabat nach außen, konstante Stoffwerte
    \cite[S.~838, Gl.~31.57--31.59]{almbauerThermomanagement2019},
    \cite[S.~256, Gl.~12.7]{pucherLadeluftkuehlungUndLadeluftkuehler2012}.
  \item $\Delta T_{\log}$ gehört zur räumlichen Beschreibung und ist nicht
    $\Tout - \TLT$. Die Gegenstromform ist idealisiert, die Bauform folgt nicht
    aus ihr; für den unten genutzten Grenzfall $C_r \to 0$ ist die Stromführung
    ohne Einfluss \cite[S.~3258 und 3264]{yinAnalyticModelTransient2003},
    \cite[Gl.~31.64]{almbauerThermomanagement2019}.
  \item $\TLT$ als gemessene Randbedingung führen, nicht als Konstante: der
    geregelte LT-Vorlauf folgt Last und Ventilhub schwach
    (Abschnitt~\ref{sec.erg.ol}). $\THT$ analog für das HT-Bündel.
  \item Stationäre Bilanz und dynamische Wassererwärmung trennen; die Wasserseite
    liefert hier nur die Randbedingung.
\end{itemize}

\begin{equation}
  \label{eq.energiebilanz}
  \dot m_L\,c_{p,L}\,(\Tin - \Tout)
  \;=\; k\,A\,\Delta T_{\log}
  \;=\; \dot m_W\,c_{p,W}\,(T_{\mathrm{LT,out}} - \TLT)
\end{equation}

% ---- Absatz 3: Kennzahlen und stationaerer Zielwert ----
\begin{itemize}
  \item Luftseitiger Temperaturänderungsgrad
    $\varepsilon = (\Tin - \Tout)/(\Tin - \TLT)$
    \cite[S.~257, Gl.~12.8]{pucherLadeluftkuehlungUndLadeluftkuehler2012}; als
    Rekuperationsgrad bei Tschöke/Pucher 2018, S.~68, Gl.~41 (Key fehlt). Von der
    Effektivität $\varepsilon_Q = \dot Q/[C_{\min}(\Tin - \TLT)]$ unterscheiden;
    beide sind nur bei $C_{\min} = \dot m_L c_{p,L}$ gleich.
  \item $\mathrm{NTU} = kA/C_{\min}$, $C_r = C_{\min}/C_{\max}$; Gegenstrom
    $\varepsilon = [1 - e^{-\mathrm{NTU}(1-C_r)}]/[1 - C_r\,e^{-\mathrm{NTU}(1-C_r)}]$;
    Anwendung am marinen Ladeluftkühler
    \cite[Anhang~B.2, S.~25, Gl.~B.6--B.8]{suiMeanValueFirst2022}.
  \item Grenzfall $C_r \to 0$: $\varepsilon = 1 - e^{-\mathrm{NTU}}$, unabhängig
    von der Stromführung. Dass die Luftseite über den ganzen Ventilhub der
    kleinere Kapazitätsstrom ist, ist eine Annahme; Beleg in
    Abschnitt~\ref{sec.erg.ol}.
  \item Stationärer Zielwert Gl.~(\ref{eq.Tout.stat}); die Differenz
    $\Tin - \TLT$ wird als Kühlerspanne bezeichnet.
    % Kein eigenes Formelzeichen: S ist in Abschnitt ssec.grundlagen.clid die
    % Empfindlichkeitsfunktion.
  \item Zwei Bündel hintereinander ergeben dieselbe Beziehung zweimal: HT-Bündel
    mit $\THT$ auf die Zwischentemperatur $T_{\mathrm{zw}}$, LT-Bündel mit
    $\TLT$ auf $\Tout$; nur das LT-Bündel trägt das Ventil. Eigene Reduktion,
    Anschluss Pucher Gl.~12.8 je Bündel; Serienschaltung am 9L51/60DF bei
    \cite[S.~89--90, Gl.~47]{rupprechtAnalysisSimulationOptimisation2016}. Bei
    positiven NTU liegt der Zielwert zwischen $\Tin$, $\THT$ und $\TLT$.
    Sensorlage und Bündelreihenfolge sind keine verifizierte Geometrie; das
    steht als Grenze im Text.
  \item Derselbe statische Kern steht in den Mittelwertmodellen
    \cite[S.~204, Gl.~7.69]{erikssonModelingControlEngines2014},
    \cite[S.~197, Gl.~23]{theotokatosCycleMeanValue2010},
    \cite[S.~405, Gl.~6]{baldiDevelopmentCombinedMean2015},
    \cite[S.~571, Gl.~8]{mrosekParameterEstimationPhysical2009}. Holmgren 2005,
    Gl.~3.26 trägt einen Druckfehler (Plus statt Minus), beim Zitieren die
    Minusform nennen (Key fehlt).
\end{itemize}

\begin{equation}
  \label{eq.Tout.stat}
  \Tout \;=\; \Tin - \varepsilon\,(\Tin - \TLT),
  \qquad \varepsilon = 1 - e^{-\mathrm{NTU}} \quad \text{für } C_r \to 0
\end{equation}
\[
  T_{\mathrm{zw}} = \THT + (\Tin - \THT)\,e^{-\mathrm{NTU}_{\mathrm{HT}}},
  \qquad
  \Tout = \TLT + (T_{\mathrm{zw}} - \TLT)\,e^{-\mathrm{NTU}_{\mathrm{LT}}(\valveCA)}
\]

% COPILOT-STICHPUNKTE BEGIN Raeumliche-Abkuehlung-20260922
\begin{itemize}
  \item \textbf{Copilot: Einordnung der weiterführenden NTU-Beziehung.}
    Die oben verwendete Exponentialform entsteht aus der räumlichen
    Bilanz, etwa
    $C_L\,\mathrm dT/\mathrm da=-k(T-T_W)$ bei entlang der Fläche
    konstantem $T_W$ \cite[S.~840, Gl.~31.64]{almbauerThermomanagement2019}.
    Sie beschreibt noch keinen zeitlichen Speicher.
    Beim Ersatz der Wand durch eine konstante Wassertemperatur
    zusätzlich $C_W\gg C_L$ voraussetzen. $C_r\to0$ allein sagt
    nicht, welche Seite den kleineren Wärmekapazitätsstrom hat.
    Allgemein gilt für den luftseitigen Temperaturänderungsgrad
    $\varepsilon=(C_{\min}/C_L)\varepsilon_Q$.
\end{itemize}
% COPILOT-STICHPUNKTE END Raeumliche-Abkuehlung-20260922

% ---- Absatz 4: Waermedurchgang, Massenstrom, Stellkanal ----
\begin{itemize}
  \item Serienwiderstand Gl.~(\ref{eq.k.zusammensetzung}): Luftübergang,
    Wandleitung, Wasserübergang
    \cite[S.~256--257, Gl.~12.6]{pucherLadeluftkuehlungUndLadeluftkuehler2012},
    \cite[S.~838, Gl.~31.59]{almbauerThermomanagement2019}; $\eta_f$ ist der
    Oberflächenwirkungsgrad der berippten Luftseite. Dominanz aus den
    Widerständen beurteilen (Rippenfläche), nicht allein aus
    $\alpha_L < \alpha_W$.
  \item Massenstromabhängigkeit: $\alpha \propto \dot m^{p}$; $p$ gehört zu einer
    konkreten Korrelation mit Geometrie- und Reynolds-Gültigkeitsbereich,
    turbulente Rohrströmung $p \approx 0{,}8$
    \cite[S.~825--840]{almbauerThermomanagement2019}. Primärbeleg für variable
    Koeffizienten am Ladeluftkühler
    \cite[Abschn.~4.3, S.~8, Gl.~20--22]{vagapovDynamicModelTemperature2022}.
    Pucher Gl.~12.7 belegt kein Potenzgesetz.
  \item Bedingte eigene Folgerung: bei luftseitiger Dominanz
    $kA \propto \dot m_L^{p}$ und $\mathrm{NTU} \propto \dot m_L^{p-1}$, also
    fällt $\varepsilon$ mit dem Luftmassenstrom
    \cite[S.~206]{erikssonModelingControlEngines2014},
    \cite[Gl.~15]{theotokatosCycleMeanValue2010};
    \cite[Gl.~B.8]{suiMeanValueFirst2022} setzt
    $\mathrm{NTU} \propto \dot m^{-0{,}2}$. Der am Prüfstand angepasste Exponent
    ist eine Fitgröße auf einem Massenproxy, kein Messwert des
    Wärmeübergangsgesetzes (Abschnitt~\ref{sec.erg.ol}).
  \item $\dot m_L$ wird am Großmotor nicht gemessen (Llamas 2019, S.~18, Key
    fehlt). Ersatz: Last, Laderdrehzahl und Ladedruck über den Liefergrad
    (Heywood, S.~54, Gl.~2.27a, Key prüfen) oder der Druckverlust über den
    Kühler, $\Delta p \propto \dot m^{2}/\rho$, also
    $\dot m \propto \sqrt{\Delta p\,p/T}$ (Beleg nachschlagen). Jeder Last- oder
    Drehzahlterm ist damit ein Massenstromterm.
  % COPILOT-STICHPUNKTE BEGIN Luftmassenstromproxy-20260922
  % Kapitelort: 2.2, nach Massenstromabhaengigkeit, vor Stellkanal.
  % Originalpunkt oben erhalten; die folgenden Punkte praezisieren ihn.
  % Datenbeleg: 05_Matlab_TB20/10_Voranalyse/
  % Luftmassenstrom_Waermeuebertragung/ERGEBNIS.md, Z. 145--168.
  % Messstellen-Grenze: 05_Matlab_TB20/10_Voranalyse/
  % Druckdifferenz/README.md, Z. 34--46.
  \item \textbf{Copilot: Ergänzung zur Herleitung des Luftmassenstromproxys
    (22.09.2026).} Präzisierung zum vorherigen Punkt: Last und Laderdrehzahl
    können mit dem Luftmassenstrom zusammenhängen; ein entsprechender
    Modellterm bildet diesen Zusammenhang nicht zwangsläufig ab. Für die
    untersuchte Datenbasis einen direkt gemessenen Luftmassenstrom von
    berechneten Referenzkanälen unterscheiden (Kapitel~\ref{ch.datenbasis}).
    Gesucht: ein aus Messgrößen berechneter Ersatz für den relativen
    Luftmassenstrom durch den HP-Kühler. Physikalischer Anschluss:
    Wärmekapazitätsstrom $C_L=\dot m_L c_{p,L}$ in
    Gl.~(\ref{eq.energiebilanz}) und massenstromabhängiger Wärmeübergang.
  \item \textbf{Druckverlust als Ausgangspunkt.} Für eine feste Geometrie,
    positive Strömungsrichtung und näherungsweise quasistationäre,
    turbulente Durchströmung bei überwiegenden Reibungsverlusten:
    quadratischer Verlustansatz mit einer
    repräsentativen Dichte $\rho_L$. Druckverlustbeiwert $\zeta_L$ fasst die
    Strömungswiderstände zusammen; seine näherungsweise Konstanz im
    betrachteten Bereich ist eine zusätzliche Annahme
    \cite[Abschnitte \emph{Momentum Balance} und \emph{Friction}]{mathworksHeatExchangerInterfaceG}.
    Eigene Zusammenfassung und algebraische Umstellung mit der
    Strömungsfläche $A_{\mathrm{Str}}$ und der Geschwindigkeit $v_L$:
    \begin{equation}
      \label{eq.luftproxy.druckverlust}
      \begin{aligned}
        \dpHP &\approx \frac{\zeta_L}{2}\rho_L v_L^2,
        \qquad \dot m_L=\rho_L A_{\mathrm{Str}}v_L,\\
        \Rightarrow\quad \dot m_L &\approx A_{\mathrm{Str}}
          \sqrt{\frac{2\,\dpHP\,\rho_L}{\zeta_L}}.
      \end{aligned}
    \end{equation}
    $A_{\mathrm{Str}}$: durchströmter Querschnitt; von der
    Wärmeübertragungsfläche $A$ in Gl.~(\ref{eq.energiebilanz}) unterscheiden.
  \item \textbf{Dichte und Temperaturwahl.} Idealgasnäherung
    $\rho_L=p_{\mathrm{abs}}/(R_L T_{\mathrm{K}})$ mit spezifischer
    Gaskonstante $R_L$ und absoluter Temperatur $T_{\mathrm{K}}$.
    Für einen aus Eintrittssignalen gebildeten Proxy: Eintrittsdichte als
    repräsentative Dichte ansetzen. Daraus folgt bei konstantem
    $\zeta_L$ und $R_L$
    \begin{equation}
      \label{eq.luftproxy.dichte}
      \dot m_L\propto
      \sqrt{\frac{\dpHP\,p_{\mathrm{col,in,abs}}}
                       {T_{\mathrm{col,in,K}}}}.
    \end{equation}
    $T_{\mathrm{col,in,K}}$: gemessenes $\Tin$ als absolute Temperatur.
    Die Luft kühlt entlang des Kühlers ab; die Eintrittsdichte ersetzt
    hier den räumlichen Dichteverlauf. Ein anderer repräsentativer
    Zustand verlangt eine dazu passende Druck- und Temperaturzuordnung.
    $\Tout$ im Nenner ist deshalb eine gesondert zu begründende Variante.
  \item \textbf{Dimensionslose Form.} Bezugsgrößen an einem gemeinsamen,
    festen Referenzzustand definieren; positive Druckdifferenz,
    absoluten Druck und absolute Temperatur verwenden:
    \begin{equation}
      \label{eq.luftproxy.relativ}
      \mrel =
      \sqrt{\frac{\dpHP}{\Delta p_{\mathrm{HP,ref}}}
            \frac{p_{\mathrm{col,in,abs}}}{p_{\mathrm{col,in,abs,ref}}}
            \frac{T_{\mathrm{col,in,K,ref}}}{T_{\mathrm{col,in,K}}}}
      \;\approx\;\frac{\dot m_L}{\dot m_{L,\mathrm{ref}}}.
    \end{equation}
    Die Näherung gilt unter den zuvor genannten Annahmen. Der
    Referenzmassenstrom muss für die Berechnung des Proxys nicht bekannt
    sein; aus den Messgrößen entsteht keine kalibrierte Angabe in
    $\mathrm{kg/s}$. Für Gas- und Dieselbetrieb dieselben Bezugsgrößen
    verwenden; ihre Festlegung gehört in Kapitel~\ref{ch.methodik}.
  \item \textbf{Vereinfachung für datengetriebene Modelle.} Feste
    multiplikative Faktoren lassen sich bei linearer oder
    potenzförmiger Verwendung des Proxys algebraisch in einen angepassten
    Modellkoeffizienten aufnehmen. Bei anderen Modellfunktionen die
    Umparametrierung gesondert prüfen. Die Normierung kann dann bei der
    Online-Auswertung entfallen; das unnormierte Merkmal ist weiterhin
    dimensionsbehaftet. Temperaturabhängigkeit und Wurzelbildung lassen
    sich im Allgemeinen nicht durch einen konstanten Faktor ersetzen.
    Die Variante ohne Temperaturkorrektur setzt stattdessen
    \[
      q_{\mathrm{rel}}=
      \sqrt{\frac{\dpHP}{\Delta p_{\mathrm{HP,ref}}}
            \frac{p_{\mathrm{col,in,abs}}}{p_{\mathrm{col,in,abs,ref}}}}
      =\mrel\sqrt{\frac{T_{\mathrm{col,in,K}}}
                            {T_{\mathrm{col,in,K,ref}}}}.
    \]
    Temperaturumrechnung in Kelvin und Umrechnung eines Relativdrucks
    auf Absolutdruck sind additive Umrechnungen. Innerhalb dieser
    Wurzelfunktion können sie nicht einfach in einem konstanten
    Koeffizienten aufgehen. Einheitenfaktoren dagegen sind konstant.
  \item \textbf{Geltungsbereich und Einordnung.} Die Herleitung enthält
    keine brennstoffspezifische Größe. Ein gemeinsamer Proxy für Gas und
    Diesel ist bei gleichem Luftpfad und vergleichbaren Stoffeigenschaften
    begründbar; gleiche Modellkoeffizienten oder gleiche Güte folgen
    daraus nicht. Reynolds-Abhängigkeit des Verlustbeiwerts,
    Kondensation, Verschmutzung und Änderungen der Strömungsaufteilung
    begrenzen die Näherung. Ein aus der Motorfüllung abgeleiteter
    Massenstrom ist bei Abzweigungen oder Speicherwirkung nicht ohne
    Bilanzprüfung dem Kühlerdurchsatz gleichzusetzen.
    Die gemessene Druckdifferenz muss den betrachteten Kühlerpfad erfassen.
    Konkrete Messkanäle, Druckabgriffe und Druckrekonstruktion:
    Kapitel~\ref{ch.datenbasis} und \ref{ch.methodik}. Vergleich von
    $\mrel$ und $q_{\mathrm{rel}}$ hinsichtlich Vorhersagefehler und
    Ventilwirkung: Kapitel~\ref{ch.ergebnisse}; daraus allein folgt kein
    Nachweis der Genauigkeit gegenüber dem tatsächlichen Luftmassenstrom.
  % COPILOT-STICHPUNKTE END Luftmassenstromproxy-20260922
  \item Stellkanal: das Ventil stellt $\dot m_W$ durch das LT-Bündel, damit
    $\alpha_W$, $kA$, NTU und $\varepsilon$. Bei großem Wasserdurchsatz
    verschwindet $1/(\alpha_W A_W)$ gegen den Luftwiderstand, weiterer Hub wirkt
    weniger. Totzone, Wendepunkt und S-Form folgen nicht aus der Bilanz; sie
    sind Messbefund (Abschnitt~\ref{ssec.erg.ol.kennlinie}).
  \item Ungemessene Hypothesen für kleine Hübe: laminarer Bodenwert der
    Nusselt-Zahl bei kleinem Wasserdurchsatz (Vagapov 2024, Gl.~7.31) und
    Wechsel des kleineren Kapazitätsstroms auf die Wasserseite (Tschöke/Pantow
    2018, S.~696, Gl.~8 und 9); beide Keys liegen in den BIB-Nachträgen.
  \item Die Ventilkennlinie selbst ist je Pfad linear (Abschnitt~\ref{ChATCo});
    $\dot m_W(\valveCA)$ und $\varepsilon(\valveCA)$ sind es nicht.
    $\varepsilon$ ist eine Fläche über Hub und Luftmassenstrom
    \cite[S.~205, Gl.~7.70 und 7.71]{erikssonModelingControlEngines2014},
    \cite[S.~132]{erikssonMODELINGTURBOCHARGEDSI},
    \cite[S.~11, Gl.~2.16]{petterssonSimulationTurboCharged2000},
    \cite[Gl.~B.8]{suiMeanValueFirst2022},
    \cite[S.~89]{rupprechtAnalysisSimulationOptimisation2016}.
\end{itemize}

\begin{equation}
  \label{eq.k.zusammensetzung}
  \frac{1}{k\,A} \;=\; \frac{1}{\eta_f\,\alpha_L\,A_L}
  + R_{\mathrm{Wand}}
  + \frac{1}{\alpha_W\,A_W}
\end{equation}

% COPILOT-STICHPUNKTE BEGIN Widerstaende-und-Stroemung-20260922
\begin{itemize}
  \item \textbf{Copilot: Grenze der Massenstrom-Potenznäherung.}
    Aus einer Nusselt-Korrelation kann bei festen Stoffwerten und
    fester Geometrie $\alpha\propto\dot m^b$ folgen.
    Bei dominierendem Luftwiderstand und $C_{\min}=C_L$ ergibt sich
    $\mathrm{NTU}\propto\dot m_L^{b-1}$; eine abnehmende
    Effektivität setzt zusätzlich $b<1$ voraus
    \cite[S.~826--827 und 839]{almbauerThermomanagement2019}.
    Diese Annahmen begründen keinen universellen Exponenten für die
    spätere datengetriebene Modellierung.
\end{itemize}
% COPILOT-STICHPUNKTE END Widerstaende-und-Stroemung-20260922

% ---- Absatz 5: Stellverstaerkung ----
\begin{itemize}
  \item Kettenregel auf Gl.~(\ref{eq.Tout.stat}) ergibt
    Gl.~(\ref{eq.grundlagen.gain.epsntu}): Spannenterm mit
    $\mathrm{d}\varepsilon/\mathrm{d}\valveCA$, dazu die Mitbewegung von $\Tin$
    (Gewicht $1-\varepsilon$) und $\TLT$ (Gewicht $\varepsilon$). Das ist
    Algebra, keine Messung; dass die beiden hinteren Terme nicht null sind,
    zeigt Abschnitt~\ref{sec.erg.ol}.
  \item Lokales Vorzeichen negativ nur unter Festhaltebedingung und bei mit dem
    Hub wachsendem $\varepsilon$. Betrag proportional zur Kühlerspanne, steigt
    also mit Last und Ladedruck; fällt auf den kleinen Restbeitrag des dritten
    Terms, sobald $\varepsilon$ nahe eins liegt (Ventil weit offen, Feldbetrieb).
    Volle Ventilautorität ist deshalb nur im Flankenbereich der Kennlinie
    messbar.
  \item Über $\mathrm{NTU}(\dot m_L)$ hängt $\partial\varepsilon/\partial\valveCA$
    selbst von der Last ab; Kennlinienform und Spanne trennen sich nur
    näherungsweise.
  \item Konvention für jede Verstärkungsangabe: Steigungsdefinition (Sekante über
    ein Hubband oder lokale Steigung), Ventilband und Arbeitspunkt. Eine
    Vollhub-Sekante ist keine Untergrenze aller lokalen Steigungen, schon wegen
    der Totzone. Kapitel~\ref{ch.ergebnisse} hält die Konvention ein.
\end{itemize}

\begin{equation}
  \label{eq.grundlagen.gain.epsntu}
  \frac{\mathrm{d}\Tout}{\mathrm{d}\valveCA}
  \;=\; -\,(\Tin - \TLT)\,\frac{\mathrm{d}\varepsilon}{\mathrm{d}\valveCA}
  \;+\; (1-\varepsilon)\,\frac{\mathrm{d}\Tin}{\mathrm{d}\valveCA}
  \;+\; \varepsilon\,\frac{\mathrm{d}\TLT}{\mathrm{d}\valveCA}
\end{equation}

% COPILOT-STICHPUNKTE BEGIN Partielle-Ventilwirkung-20260922
\begin{itemize}
  \item \textbf{Copilot: Festhaltebedingung der Stellwirkung.}
    Aus Gl.~(\ref{eq.grundlagen.basic.stufen}) folgt bei unverändertem
    $T_{\mathrm{zw}}$, $\TLT$ und Luftmassenstrom
    $\partial\Tstat/\partial\valveCA
      =-(T_{\mathrm{zw}}-\TLT)
         \partial\varepsilon_{\mathrm{LT}}/\partial\valveCA$.
    Das Vorzeichen hängt damit auch von der hydraulischen Stellrichtung
    und der Temperaturdifferenz ab. Die Gesamtmitbewegung im geregelten
    Betrieb ist eine andere Größe als diese partielle Stellwirkung.
\end{itemize}
% COPILOT-STICHPUNKTE END Partielle-Ventilwirkung-20260922

% ---- Absatz 6: Speicher, Rate, Totzeit ----
\begin{itemize}
  \item Wand und Kühlwasser sind eigene Speicher
    \cite[S.~830--831, Gl.~31.35--31.38; S.~838--840, Gl.~31.57--31.65]{almbauerThermomanagement2019},
    \cite[Gl.~1 und 2]{yinAnalyticModelTransient2003},
    \cite[Gl.~2 und 27]{ariciVehicleEngineCooling1999}. Austrittstemperatur und
    Wandtemperatur nicht gleichsetzen.
  \item Eigene Reduktion auf einen Zustand, Gl.~(\ref{eq.energiebilanz.dyn}):
    konvektiver Eintrag aus der Behälterbilanz
    \cite[S.~31, Gl.~2.7]{guzzellaIntroductionModelingControl2010},
    \cite[S.~164, Gl.~7.26a]{erikssonModelingControlEngines2014},
    \cite[S.~147, Gl.~4.1.48]{isermannEngineModelingControl2014} minus
    Wärmeabfuhr. $C_{\mathrm{eff}}$ bündelt die Speicher; $U_{\mathrm{eff}}A$
    ist ein effektiver Leitwert, nicht das $kA$ der Statik. Vorbild für eine
    vergrößerte Speichermasse ohne Zusatzzustand
    \cite[S.~579, Gl.~26]{mrosekParameterEstimationPhysical2009}.
  \item Zeitkonstante bei eingefrorenen Parametern
    $\tau = C_{\mathrm{eff}}/(\dot m_L c_{p,L} + U_{\mathrm{eff}}A)$ (Stanivuk
    2021, S.~117, Gl.~9, Key fehlt;
    \cite[Gl.~31.36]{almbauerThermomanagement2019}). Wird die Mitteltemperatur
    $(\Tin + \Tout)/2$ als Treiber angesetzt, steht $kA/2$ im Nenner
    \cite[S.~230, Gl.~4.6.14--4.6.16]{isermannEngineModelingControl2014}; die
    Größenordnung ändert das nicht. Pole in 1/s, Zeitkonstanten in s.
  \item Folge für die Rate: sie wächst mit dem Luftmassenstrom und über
    $kA(\dot m_W)$ schwach mit dem Wasserdurchsatz
    \cite[S.~248, Gl.~4.7.39]{isermannEngineModelingControl2014}. Größenordnung
    nicht aus Auslegungsdaten festlegen, sondern aus Daten
    (Kapitel~\ref{ch.methodik} und \ref{ch.ergebnisse}); die Autokorrelation
    geregelter Betriebsdaten identifiziert weder den offenen Pol noch dessen
    Speicher.
  \item Konsistenzgrenze: stationär liefert die Einzustandsform
    $\varepsilon = N/(1+N)$ mit $N = U_{\mathrm{eff}}A/(\dot m_L c_{p,L})$, nicht
    $1 - e^{-\mathrm{NTU}}$. Deshalb Bilanz zur Termmotivation,
    $\varepsilon$-NTU-Statik zum Zielwert: Trennung in nichtlineare Statik und
    nachgeschaltete Dynamik, Hammerstein-Form
    \cite[Abschn.~2, Abb.~2, S.~4]{vagapovDynamicModelTemperature2022},
    \[
      \dot y = \lambda(\dot m_L, \valveCA)\,(T_{\mathrm{stat}} - y)
    \]
    mit $T_{\mathrm{stat}}$ aus Gl.~(\ref{eq.Tout.stat}). Bei veränderlicher
    Rate ist die Dynamik zeitvariant; ein zeitinvariantes lineares Glied wäre
    eine Zusatzannahme.
  \item Modellordnung ist Prüfaufgabe, keine feste Zahl: bei Vagapov zwei
    dominante Beiträge der Kühlmedien, Reduktion auf einen dort mit schlechterem
    Verlauf \cite[Abschn.~5.2, Gl.~43]{vagapovDynamicModelTemperature2022}; der
    zweite Pol liegt auf der Kaltseite (Vagapov 2024, Abb.~7.10, Dominanzmaß).
    Prüfung für die eigene Anlage in Kapitel~\ref{ch.methodik} und
    \ref{ch.ergebnisse}.
  \item Totzeit als Verweilzeit $m_s/\dot m_L$
    \cite[Gl.~41]{vagapovDynamicModelTemperature2022},
    \cite[S.~3257]{yinAnalyticModelTransient2003},
    \cite[Gl.~31.57]{almbauerThermomanagement2019},
    \cite[Gl.~2.173]{guzzellaIntroductionModelingControl2010} plus Rohrstrecke
    zwischen Ventil und Kühler (Abschnitt~\ref{ChATCo}). Messkette als
    Sensor-PT1 \cite[Gl.~4.6.42]{isermannEngineModelingControl2014},
    \cite[Tab.~7.6]{erikssonModelingControlEngines2014},
    \cite[Gl.~17]{mrosekParameterEstimationPhysical2009}.
\end{itemize}

\begin{equation}
  \label{eq.energiebilanz.dyn}
  C_{\mathrm{eff}}\,\frac{\mathrm{d}\Tout}{\mathrm{d}t}
  \;=\;
  \underbrace{\dot m_L\,c_{p,L}\,(\Tin - \Tout)}_{\text{Lufteintrag}}
  \;-\;
  \underbrace{U_{\mathrm{eff}}A\,(\Tout - \TLT)}_{\text{Wärmeabfuhr}}
\end{equation}

% COPILOT-STICHPUNKTE BEGIN Speicherherleitung-20260922
\begin{itemize}
  \item \textbf{Copilot: Grenze der Einzustandsbilanz.}
    Gl.~(\ref{eq.energiebilanz.dyn}) beschreibt einen einstufigen
    Ersatzkühler mit $\Tin$ und $\TLT$. Der unabhängige HT-Beitrag
    entfällt durch eine zusätzliche Modellannahme.
    Für die LT-Stufe allein wäre $T_{\mathrm{zw}}$ die Eintrittstemperatur.
    Ihr stationärer Grenzwert und der Zielwert der HT-/LT-Exponentialform
    sind im Allgemeinen verschieden. Eine Relaxation zum
    HT-/LT-Zielwert gemäß Gl.~(\ref{eq.grundlagen.basic.dynamik})
    ist deshalb eine eigene reduzierte Modellform.
  \item Die Verteilung des Ventileinflusses auf Zielwert und Rate sowie
    seine Stärke sind Modellentscheidungen. Thermische Speicherung,
    Transportzeit und Sensordynamik können zur gemessenen Verzögerung
    beitragen \cite[S.~230 und 237--238]{isermannEngineModelingControl2014};
    eine identifizierte Zeitkonstante ist keinem Bauteil eindeutig zugeordnet.
\end{itemize}
% COPILOT-STICHPUNKTE END Speicherherleitung-20260922

\anno{Bild (Visio): thermisches Ersatzschaltbild, $C_{\mathrm{eff}}$ als
Kapazität, $1/(kA)$ als Widerstand mit Ventilpfeil auf den wasserseitigen
Anteil, $\Tin$ und $\TLT$ als Quellen, HT-Bündel als vorgeschalteter Block.}

% ---- Absatz 7: lokale Naeherung und Produktfamilien ----
\begin{itemize}
  \item Mit $d = \Tin - \Tout$ und der Kühlerspanne gilt
    $\Tout - \TLT = (\Tin - \TLT) - d$; Gl.~(\ref{eq.energiebilanz.dyn}) wird zu
    \[
      \frac{\mathrm{d}\Tout}{\mathrm{d}t} = a\,d - b\,(\Tin - \TLT),
      \qquad
      a = \frac{\dot m_L c_{p,L} + U_{\mathrm{eff}}A}{C_{\mathrm{eff}}},
      \quad
      b = \frac{U_{\mathrm{eff}}A}{C_{\mathrm{eff}}}.
    \]
    Eigene Umformung, keine Literaturgleichung.
  \item $a$ und $b$ hängen vom nicht gemessenen $\dot m_L$ ab. Lokal affin in
    messbaren Betriebspunktgrößen (Drehzahl-, Last-, Spannenkoordinate) ergeben
    sich zwei Produktfamilien: Proxy mal $d$ und Proxy mal Spanne. Die Physik
    begründet die Familien; Proxywahl und affine Näherung sind
    Modellentscheidungen (Kapitel~\ref{ch.methodik}). Stütze:
    \cite[S.~132]{erikssonMODELINGTURBOCHARGEDSI}, Effektivität abhängig von
    Temperaturniveau, $\dot m_L$ und $\dot m_L/\dot m_W$ (Luft-Luft-Kühler,
    Struktur übertragbar, Koeffizienten nicht).
  \item Ventil und Temperatur koppeln in der Bilanz über
    $kA(\dot m_W)\,(\Tout - \TLT)$. In der Statikform steckt diese Kopplung in
    $\varepsilon(\valveCA)\,(\Tin - \TLT)$: eine Funktion des Hubs mal exogene
    Temperaturen, nicht mal $\Tout$. Das ist der physikalische Grund, den
    Ventileinfluss über den Zielwert zu führen; die Rate trägt nur den schwachen
    Rest aus dem vorigen Absatz.
  \item Konstante und linearer $\Tout$-Term sind als affine Restkorrektur
    zulässig, aber keinem Wärmestrom und keinem Speicher zuzuordnen; ein
    zusätzlicher $\Tout$-Term verschiebt zusammen mit dem Kern den Gesamtpol.
  \item Eingangsvorgeschichte: eine kausal gefilterte Spanne kann nicht
    aufgelöste Speicher (Wand, Wasser) vertreten; als Filterzustand exakt
    definiert, physikalisch nur Hypothese
    \cite[Abschn.~5, S.~10--12]{vagapovDynamicModelTemperature2022}.
\end{itemize}

% ---- Absatz 8: was kein Bibliotheksterm wird ----
\begin{itemize}
  \item Umgebungsverlust $(\Tout - T_{\mathrm{amb}})$
    \cite[Gl.~27 und 34]{ariciVehicleEngineCooling1999},
    \cite[Gl.~31.13]{almbauerThermomanagement2019},
    \cite[Gl.~10]{vagapovDynamicModelTemperature2022}: die adiabate Bilanz ist
    Annahme, Prüfung in Abschnitt~\ref{sec.erg.ol}; $T_{\mathrm{amb}}$ bleibt
    Gültigkeitsgrenze.
  \item Kondensation
    \cite[Abschn.~12.4, S.~265--268]{pucherLadeluftkuehlungUndLadeluftkuehler2012};
    der marine Kühler wird bei Arava 2026 ohne latenten Term gerechnet (Key
    fehlt). Kein Term (Feuchte nicht erfasst, unstetig), sondern Grenze:
    Taupunkt am Kapitelanfang, weiche Untergrenze des Reglers in
    Kapitel~\ref{ch.ergebnisse}.
  \item Brennstoffmodus kommt in der Bilanz nicht vor und wirkt nur über die
    Eingänge. Ein Niveauversatz zwischen den Modi ist deshalb kein Bilanzterm,
    sondern eine Arbeitspunkt- und Gültigkeitsfrage
    (Kapitel~\ref{ch.ergebnisse}).
    % optional, ein Satz
  \item Weitere Kandidaten aus der Literatur, hier ohne Verdikt: Füll- und
    Entleerterm $\Tout\,\dot p/p$
    \cite[Gl.~7.26a]{erikssonModelingControlEngines2014},
    \cite[Gl.~30]{theotokatosCycleMeanValue2010},
    \cite[Gl.~20.26]{merkerThermodynamischeGrundlagen2019}; Stoffwertterm über
    das Temperaturniveau $(\Tin + \TLT)/2$
    \cite[Gl.~7.70]{erikssonModelingControlEngines2014} (Taler 2017, Gl.~20,
    Key fehlt); wasserseitige Bilanz $\dot m_W\,(T_{\mathrm{LT,out}} - \TLT)$
    \cite[S.~96, Gl.~2.170 und 2.171]{guzzellaIntroductionModelingControl2010},
    \cite[Gl.~4.7.36]{isermannEngineModelingControl2014}. Prüfung in
    Kapitel~\ref{ch.methodik} und \ref{ch.ergebnisse}.
  \item Die Lücke, dass keine der Modellquellen dem Ladeluftkühler einen
    Temperaturzustand mit Ventileingang gibt (marin statisch:
    \cite{theotokatosCycleMeanValue2010}, \cite{baldiDevelopmentCombinedMean2015},
    \cite{suiMeanValueFirst2022}; Gegenposition Wahlström 2009, S.~7, Key
    fehlt), gehört nach Kapitel~\ref{ch.standderTechnik}; hier höchstens ein
    Satz.
\end{itemize}

% ---- Absatz 9: Schluss, zwei Saetze, kein Wegweiser ----
% COPILOT-STICHPUNKTE BEGIN Modellanschluss-20260922
\begin{itemize}
  \item \textbf{Copilot: Bezug für die spätere Modellierung.}
    Die einfache Herleitung am Abschnittsanfang trägt den Verweis
    auf Temperaturdifferenzen, Durchsatzprodukte und Relaxation.
    Die folgenden Detailvorschläge legen noch keine Ventilkennlinie,
    feste Modellordnung oder ausschließliche Ventilwirkung im
    stationären Zielwert fest.
\end{itemize}
% COPILOT-STICHPUNKTE END Modellanschluss-20260922
\begin{itemize}
  \item Termfamilien aus der Herleitung: stationärer Zielwert aus der seriellen
    $\varepsilon$-NTU-Beziehung mit dem Ventil nur darin; Relaxation eines
    Zustands mit massenstromabhängiger Rate; Produkte aus Betriebspunktproxy und
    Temperaturdifferenz als lokale Näherung; Totzeit und Messkette; die Grenzen
    aus dem vorigen Absatz.
  \item Auswahl, Kollinearität, Schätzbarkeit und Nutzen im Freilauf prüft
    Kapitel~\ref{ch.methodik} und \ref{ch.ergebnisse}. Ein Satz mit
    \verb|\ref|, kein Absatz.
\end{itemize}
% COPILOT-STICHPUNKTE ENDE ch2-08

% PRUEFPUNKTE (Copilot 16.09.2026, vor dem Ausformulieren):
% - Exponent: die Altfassung behauptete 0,2 aus Nu ~ Re^0,8 und "Band um 0,2".
%   Der S-Phys-Fit hat alpha = -0,53 auf m_rel (Steckbrief Physikalischer_Kern).
%   Hier nur "Fitgroesse auf dem Proxy, kein Waermeuebergangsgesetz"; Zahl und
%   Einordnung nach Kap. 6.
% - Beleg fuer den Druckverlust-Proxy (Absatz 4) fehlt: Almbauer Kap. 31 oder
%   Guzzella/Onder Kap. 2 nachschlagen.
% - Schema (Visio) nach dem Skill abbildung; Caption ein Satz.
% KEYS FEHLEN in Masterarbeit.bib (Klartext oben durch \cite ersetzen):
% - Staging BIB_NACHTRAG_0709.bib: Holmgren 2005, Stanivuk 2021, Tschoeke/Pucher
%   2018 (+ Kapitel Pucher, Pantow), Ekberg 2018, Arava 2026, Taler 2017.
% - Staging BIB_NACHTRAG_2608.bib: Vagapov 2024 (vagapovModellierungIdentifikationUnd2024).
% - Heywood, Llamas 2019, Wahlstroem 2009: Key pruefen.
% GESTRICHEN gegenueber der Altfassung: eq.verdichter, eq.alpha.scaling,
% eq.kA.Praxis (nur ch2-intern referenziert, jetzt Saetze); Tabelle
% tab.grundlagen.terme.literatur samt \anno; alle eigenen Pruefstandszahlen
% (Saettigung, Totzone, epsilon-Band, Zeichnungsabschaetzung, NTU-Band);
% Copilot-Zusaetze ch2-05/06/07 und Kommentarbloecke vom 07.09. (aufgegangen).






%%%%%%%%%%%%%%%%%%%%%%%%%%--------------------------------------------
%  Modellprädiktive Regelung
%%%%%%%%%%%%%%%%%%%%%%%%%%--------------------------------------------

\section{Modellprädiktive Regelung}
Heute regelt die Anlage mit dem PI-Anteil und der Vorsteuerung aus
Abschnitt~\ref{ChATCo}. Der PI-Anteil reagiert auf die
aktuelle Regelabweichung und greift damit erst ein, wenn die Ladelufttemperatur
bereits abweicht. Der Stellkanal ist dabei nichtlinear, mit Totzone und Sättigung
in der Ventilkennlinie (Abschnitt~\ref{ssec.erg.ol.kennlinie}). Großmotoren
dieser Baugröße sind zudem träge, sodass eine Abweichung spät zurückgeholt wird.
Ein modellprädiktiver Regler rechnet den Verlauf dagegen über einen Horizont
voraus und führt Stell- und Temperaturgrenzen als Nebenbedingung mit.

% ===== ERGÄNZT Copilot 01.09.2026 (PI-Gesetz, MPC-Gleichung, Offset-frei;
% ===== Fassung nach Lektor: K_c/T_I-Form mit Vorsteueranteil, weiche
% ===== Untergrenze, Stellhorizont N_u, u_{-1} definiert) =====
Der Serienregler kombiniert eine Vorsteuerung mit einem PI-Anteil
(Abschnitt~\ref{ChATCo}). Die Stellgröße folgt aus der Regelabweichung
$e(t) = r(t) - y(t)$,
\begin{equation}
  \label{eq.grundlagen.pi}
  u(t) \;=\; u_{\mathrm{ff}}(t)
  + K_c \Bigl( e(t) + \frac{1}{T_I} \int_0^t e(\tau)\,\mathrm d\tau \Bigr),
\end{equation}
mit dem Vorsteueranteil $u_{\mathrm{ff}}$ aus dem Kennfeld über Last
und Temperaturdifferenz, der Reglerverstärkung $K_c$ und der
Nachstellzeit $T_I$. Der Integralanteil beseitigt bleibende
Abweichungen und trägt damit das Niveau; er greift aber erst ein,
nachdem sich eine Störung im Ausgang gezeigt hat. An den Stellgrenzen
hält ein Anti-Windup den Integrator an (Abschnitt~\ref{ChATCo}).

Der modellprädiktive Regler löst dagegen in jedem Takt ein
Optimierungsproblem über einen Prädiktionshorizont von $N$ Schritten
mit $N_u \le N$ freien Stellschritten,
\begin{equation}
  \label{eq.grundlagen.mpc}
  \min_{\Delta u_0,\dots,\Delta u_{N_u-1}}\;
  \sum_{k=1}^{N} \bigl(y_k - r_k\bigr)^2
  \;+\; \lambda \sum_{k=0}^{N_u-1} \Delta u_k^2
  \;+\; \mu \sum_{k=1}^{N} s_k^2,
\end{equation}
unter $u_{\min} \le u_k \le u_{\max}$ und der weichen Untergrenze
$y_k \ge y_{\min} - s_k$ mit Schlupfvariablen $s_k \ge 0$. Die Folge
$y_k$ liefert das Streckenmodell; $\Delta u_k = u_k - u_{k-1}$ ist der
Stellschritt, wobei $u_{-1}$ die zuletzt gestellte Größe bezeichnet.
Das Gewicht $\lambda$ setzt Stellbewegung ins Verhältnis zum
Regelfehler, $\mu$ bestraft Bandverletzungen
\cite[Kap.~1]{rawlingsModelPredictiveControl20202020}. Angewendet wird
nur der erste Stellschritt; im nächsten Takt wird mit der neuen Messung
neu gerechnet (gleitender Horizont). Für die Ladeluftstrecke ist $u$
die Ventilstellung $\valveCA$. Ist das Modell linear in der
Stellgröße, ist (\ref{eq.grundlagen.mpc}) ein konvexes quadratisches
Programm. Bleibende Störungen und Modellfehler behandelt die
offsetfreie Erweiterung: Ein Störmodell wird als zusätzlicher Zustand
mitgeschätzt, ein Beobachter korrigiert damit die Vorhersage, und der
Regler erreicht den Sollwert trotz Versatz
\cite{muskeDisturbanceModelingOffsetfree2002,
pannocchiaDisturbanceModelsOffsetfree2003,
maederLinearOffsetfreeModel2009}. Für die Umsetzung wird das
Streckenmodell mit einem Halteglied im Reglertakt diskretisiert; Takt,
Horizont und Gewichte der vorliegenden Arbeit stehen in
Abschnitt~\ref{sec.methodik.mpc}.

Zur modellprädiktiven Regelung von Verbrennungsmotoren liegen zahlreiche Arbeiten
vor; Kapitel~\ref{ch.standderTechnik} ordnet sie ein.


\begin{figure}[htbp]
  \centering
  \includegraphics[width=\textwidth]{kapitel2/pictures/mpc_gleitender_horizont.png}
  \caption[MPC mit gleitendem Horizont]{MPC mit gleitendem Horizont:
    Temperaturvorhersage, Stellhorizont und zuerst angewandter Stellschritt.
    Eigene schematische Darstellung nach
    \cite[Kap.~1]{rawlingsModelPredictiveControl20202020};
    frei gewählte Verläufe ohne Messdaten, erneute Optimierung im nächsten Takt.}
  \label{fig.grundlagen.mpc.horizont}
\end{figure}




%Störgrößen prädiktion auf jeden falll noch ERwähnen!!!

%was ist sonst noch für mpc als grundlagenwissen für meine arbeit wichtig??